\chapter{Mạng xã hội, AI và những màn khôn lỏi online}
\markboth{Mạng xã hội, AI và những màn khôn lỏi online}{Mạng xã hội, AI và những màn khôn lỏi online}

\section{Em hỏi AI thôi chứ đâu có gian lận}
\begin{truyen}{Em hỏi AI thôi chứ đâu có gian lận}{Classroom Talk}
\chuhoa{M}{ột bạn nói rất tỉnh: ``Em chỉ nhờ AI gợi ý thôi chứ có quay bài đâu. Bài em nộp vẫn là em bấm gửi mà thầy.''}

Thầy hỏi lại: ``Nếu người khác nghĩ hộ, ráp hộ, viết hộ, còn em chỉ bấm nộp, vậy phần nào trong đó là não của em đang đi học?''

Bạn ấy cười trừ: ``Thì... em cũng có đọc qua.''

Thầy chốt: ``Đọc qua ý người khác không tự động biến nó thành hiểu của mình. Đi thuê não dài ngày thì điểm có thể qua, đầu em thì teo.''

\textit{(Using a smart tool is not the problem. Handing over your thinking and still claiming understanding is.)}
\end{truyen}

\begin{baihoc}
AI là công cụ tốt. Nhưng dùng nó để khỏi nghĩ thì chẳng khác nào đi tập gym mà thuê người khác hít đất hộ.
\textit{(AI can be a strong tool. Using it to avoid thinking is like going to the gym and paying someone else to do the push-ups.)}
\end{baihoc}

\begin{ghinhoanh}
\textbf{Short English to remember:}

Tools may assist.

They must not replace your thinking.
\end{ghinhoanh}

\begin{tuhocnhanh}
1. Em đang dùng công cụ để hiểu sâu hơn, hay để né phần khó nhất? \textit{(Do you use tools to understand better, or to dodge the hardest part?)}

2. Nếu bỏ công cụ ra, em còn tự giải thích được bao nhiêu phần? \textit{(If the tool disappears, how much can you still explain on your own?)}
\end{tuhocnhanh}
\ngancach

\section{Em đăng story than áp lực chứ em vẫn ổn}
\begin{truyen}{Em đăng story than áp lực chứ em vẫn ổn}{Classroom Talk}
\chuhoa{C}{ó bạn bảo: ``Em than trên story cho nhẹ đầu thôi. Chứ em vẫn ổn, thầy đừng làm quá.''}

Thầy hỏi: ``Ổn thật thì tốt. Nhưng nếu em chăm phát tín hiệu cầu cứu hơn chăm xử lý cái đang bóp cổ mình, vậy em đang chữa hay đang phát sóng?''

Bạn ấy nhăn mặt: ``Than chút cũng không được hả thầy?''

Thầy đáp: ``Được. Nhưng nếu status dài hơn thời gian em ngồi xử lý vấn đề, thì em đang đầu tư vào hình ảnh mệt mỏi nhiều hơn vào chuyện thoát mệt.''

\textit{(Expressing pressure is valid. Turning pressure into performance without action keeps the wound loud and unchanged.)}
\end{truyen}

\begin{baihoc}
Than đúng chỗ giúp nhẹ đầu. Than thành thương hiệu cá nhân thì vấn đề vẫn ngồi nguyên đó cười em.
\textit{(Venting in the right place can help. Turning stress into a personal brand leaves the problem exactly where it was.)}
\end{baihoc}

\begin{ghinhoanh}
\textbf{Short English to remember:}

Venting is not the same as healing.

Attention does not solve pressure.
\end{ghinhoanh}

\begin{tuhocnhanh}
1. Em đang chia sẻ để được giúp, hay để được thấy? \textit{(Do you share to get help, or just to be seen?)}

2. Việc nhỏ nào em có thể làm ngay để bớt áp lực thật? \textit{(What small action could reduce your pressure for real?)}
\end{tuhocnhanh}
\ngancach

\section{Em không flex, em chỉ cập nhật cuộc sống thôi}
\begin{truyen}{Em không flex, em chỉ cập nhật cuộc sống thôi}{Classroom Talk}
\chuhoa{M}{ột bạn phân bua: ``Em đâu có khoe. Em chỉ đăng cho vui để mọi người biết em vẫn ổn, vẫn sang, vẫn chill.''}

Thầy cười: ``Nếu một cái ổn phải đăng mười bảy góc, ba filter và bốn dòng caption để chứng minh, thường bên trong nó chưa ổn lắm đâu.''

Bạn ấy chống chế: ``Thời nay ai chẳng sống online.''

Thầy gật: ``Sống online không sai. Chỉ là đừng chỉnh mặt tiền quá kỹ tới mức quên sửa cái nền nhà đang nứt.''

\textit{(There is nothing wrong with sharing life online. The danger begins when public appearance replaces private repair.)}
\end{truyen}

\begin{baihoc}
Hình ảnh đẹp có thể câu tim vài phút. Nội lực mới giữ đời mình lâu năm.
\textit{(A polished image can win hearts for minutes. Inner substance is what carries a life for years.)}
\end{baihoc}

\begin{ghinhoanh}
\textbf{Short English to remember:}

Appearance can impress.

Substance is what sustains.
\end{ghinhoanh}

\begin{tuhocnhanh}
1. Em đang sửa đời mình, hay sửa cách đời mình nhìn từ bên ngoài? \textit{(Are you improving your life, or just improving how it looks from outside?)}

2. Nếu không đăng nữa, em còn thấy điều đó đáng sống không? \textit{(If you could not post it, would it still feel worth living?)}
\end{tuhocnhanh}
\ngancach

\section{Em tắt cam chứ đâu có tắt ý thức}
\begin{truyen}{Em tắt cam chứ đâu có tắt ý thức}{Classroom Talk}
\chuhoa{C}{ó bạn học online quen miệng: ``Em tắt cam để đỡ ngại thôi, chứ em vẫn tập trung lắm.''}

Thầy hỏi: ``Tập trung vậy thầy vừa hỏi câu gì?''

Bạn ấy im đúng kiểu mạng lag trong tâm hồn rồi nói: ``Dạ... em nghe được đoạn cuối.''

Thầy chốt: ``Ừ. Nhiều bạn tắt cam để đỡ ngại. Xong tiện tay tắt luôn trách nhiệm, phản xạ và cả sự hiện diện của mình.''

\textit{(Privacy may be valid. But in many cases, turning off the camera slowly turns off accountability too.)}
\end{truyen}

\begin{baihoc}
Không phải cứ khuất mặt là vô tội. Có kiểu mất hình rồi mất luôn cả ý thức học.
\textit{(Being unseen does not make you innocent. Sometimes disappearing visually leads to disappearing mentally as well.)}
\end{baihoc}

\begin{ghinhoanh}
\textbf{Short English to remember:}

Privacy is one thing.

Absence disguised as presence is another.
\end{ghinhoanh}

\begin{tuhocnhanh}
1. Em cần riêng tư thật, hay đang muốn lặn cho tiện? \textit{(Do you really need privacy, or are you just trying to disappear conveniently?)}

2. Khi không ai nhìn, em còn học nghiêm như lúc bị nhìn không? \textit{(When no one watches, do you study as seriously as when they do?)}
\end{tuhocnhanh}
\ngancach